\documentclass{article}
\usepackage[utf8]{inputenc}
\usepackage{amsmath}
\usepackage{amssymb}
\usepackage{stmaryrd}
\usepackage[left=2cm, right=2cm]{geometry}

\title{Homework 2 - Problem 1 \\ CS 6347}
\author{Christian Parker \\ CXP220034}

\begin{document}

\maketitle

\begin{enumerate} % Problem 1
    \item[1.] % Problem 1.1
        The smallest \(2n\)-cycle with \(n \geq 2\) is the 4-cycle graph.
        When \(0 < k < 2\), every vertex has only one color option.
        Since every vertex in the 4-cyle graph has two neighbors and k-coloring
        requires neighboring vertices to be different colors, the graph cannot be colored.
    \item[2.] % Problem 1.2
    Given the new potential function, the MAP assignment would be:
    \[
        \hat{X} = \arg\max_{X} \frac{1}{Z} \prod_{i \in V} e^{w_{x_i}} \prod_{(i,j) \in E} 1_{x_i \neq x_j}
    \]
    To determine the maximum estimate, we only need to find the valid assignment that maximizes
    the potential function \(\phi_i(x_i) = e^{w_{x_i}}\) across all vertices. Z is the same
    for all color assignments and the indicator function simply determines if the coloring
    is valid and doesn't contribute the probability beyond that. Therefore, the estimate is
    maximized when \(\prod_{i \in V} e^{w_{x_i}}\) is maximized. This function increases
    exponentially with \(w_{x_i}\) and \(w_a = a\). Therefore, we can simplify
    the problem to the following:
    \[
        \hat{X} = \arg\max_{X} \prod_{i \in V} x_i \prod_{(i,j) \in E} 1_{x_i \neq x_j}
    \]
    This MAP assignment will assign \(x_i = k\) to as many vertices as possible,
    followed by \(k-1\), and so on. The fewest colors required for a 2n-cycle graph
    is two as each vertex will be assigned one color while both of it's neighbors are
    assigned the other. So, the MAP assignment for any 2n-cycle graph would
    be the graph such that every vertex alternates between \(k\) and \(k-1\).

    \item[3.] % Problem 1.3
    The marginal probability distribution over any variable \(x_i\) is given by:
    \[
        p(x_i) = \sum_{X \backslash x_i} \Bigg( \frac{1}{Z} \prod_{i \in V} e^{x_i} \prod_{(i,j) \in E} 1_{x_i \neq x_j} \Bigg)
    \]
    For example, if \(x_1 = 1\), then the valid colorings would be:
    \[
        \begin{array}{c|c|c|c|c}
            x_1 & x_2 & x_3 & x_4 & \prod_{i \in V} e^{x_i} \\
            \hline
            1 & 2 & 1 & 2 & e^6 \\
            1 & 2 & 1 & 3 & e^7 \\
            1 & 2 & 3 & 1 & e^7 \\
            1 & 2 & 3 & 2 & e^8 \\

            1 & 3 & 1 & 2 & e^7 \\
            1 & 3 & 1 & 3 & e^8 \\
            1 & 3 & 2 & 1 & e^7 \\
            1 & 3 & 2 & 3 & e^9 \\
        \end{array}
    \]
    Therefore:
    \begin{equation}
        \begin{aligned}
            p(x_1=1) &= \sum_{X \backslash x_1} \Bigg( \frac{1}{Z} \prod_{i \in V} e^{x_i} \prod_{(i,j) \in E} 1_{x_i \neq x_j} \Bigg) \\
            &= \frac{1}{Z} \Big( e^6 + e^7 + e^7 + e^8 + e^7 + e^8 + e^7 + e^9 \Big) \\
            &= \frac{1}{Z} \Big( e^6 + 4e^7 + 2e^8 + e^9 \Big)
        \end{aligned}
    \end{equation}
    Repeating this for all possible values of \(x_i\) will give the marginal
    probability distribution over \(x_i\).
\end{enumerate}

\end{document}